% SEGMENT OLP-0082-S01
% Part: first-order-logic
% Chapter: sequent-calculus
% Section: equality

\documentclass[../../../include/open-logic-section]{subfiles}

\begin{document}

\olfileid{fol}{seq}{ide}

\olsection{\usetoken{S}{identity} உடன் \usetoken{P}{derivation}}

!!{identity} உடைய !!^{derivation}s க்குக் கூடுதல் தொடக்கத் தொடரணிகளும்
அனுமான விதிகளும் தேவை.

\begin{defn}[$\eq$ க்கான தொடக்கத் தொடரணிகள்]
$t$ ஒரு மூடிய சொல் எனில், ${} \Sequent \eq[t][t]$ ஒரு தொடக்கத் தொடரணி.
\end{defn}

$\eq$ க்கான விதிகள் பின்வருமாறு ($t_1$, $t_2$ ஆகியவை மூடிய சொற்கள்):

\begin{defish}
\Axiom$ \eq[t_1][t_2], \Gamma \fCenter \Delta, !A(t_1) $
\RightLabel{$\eq$}
\UnaryInf$\eq[t_1][t_2], \Gamma \fCenter \Delta, !A(t_2)$
\DisplayProof
\hfill
\Axiom$\eq[t_1][t_2], \Gamma \fCenter \Delta, !A(t_2) $
\RightLabel{$\eq$}
\UnaryInf$\eq[t_1][t_2], \Gamma  \fCenter \Delta, !A(t_1)$
\DisplayProof
\end{defish}

% SEGMENT OLP-0082-S02
\begin{ex}
$s$, $t$ ஆகியவை மூடிய சொற்கள் எனில், $\eq[s][t], !A(s) \Proves !A(t)$:
\begin{prooftree}
\Axiom$ !A(s) \fCenter !A(s)$
\RightLabel{\LeftR{\Weakening}}
\UnaryInf$\eq[s][t], !A(s)  \fCenter !A(s)$
\RightLabel{$\eq$}
\UnaryInf$\eq[s][t], !A(s)  \fCenter !A(t)$
\end{prooftree}
இது ஒன்றானவற்றின் பிரதியிடத்தக்கதன்மை நெறி அல்லது லைப்னிட்சின் விதி என்று
அறியப்பட்டிருக்கலாம்.

% SEGMENT OLP-0082-S03
$\Log{LK}$ ஆனது $\eq$ சமச்சீர் மற்றும் கடப்புப் பண்புடையது என்று நிறுவுகிறது:
\begin{prooftree}
\Axiom$ \fCenter \eq[t_1][t_1] $
\RightLabel{\LeftR{\Weakening}}
\UnaryInf$ \eq[t_1][t_2] \fCenter \eq[t_1][t_1] $
\RightLabel{$\eq$}
\UnaryInf$ \eq[t_1][t_2] \fCenter \eq[t_2][t_1]$
\DisplayProof\qquad\bottomAlignProof
\Axiom$ \eq[t_1][t_2] \fCenter \eq[t_1][t_2] $
\RightLabel{\LeftR{\Weakening}}
\UnaryInf$\eq[t_2][t_3], \eq[t_1][t_2]  \fCenter \eq[t_1][t_2] $
\RightLabel{$\eq$}
\UnaryInf$\eq[t_2][t_3], \eq[t_1][t_2]  \fCenter \eq[t_1][t_3]$
\RightLabel{\LeftR{\Exchange}}
\UnaryInf$\eq[t_1][t_2], \eq[t_2][t_3]  \fCenter \eq[t_1][t_3]$
\end{prooftree}
இடப்புற !!{derivation}[loc], நமது $!A(x)$ என்பது !!{formula}~$\eq[x][t_1]$.
வலப்புறம் $!A(x)$ ஐ~$\eq[t_1][x]$ என்று கொள்கிறோம்.
\end{ex}

% SEGMENT OLP-0082-S04
\begin{prob}
பின்வரும் தொடரணிகளுக்கு !!{derivation}s தருக:
\begin{enumerate}
\item $\Sequent \lforall[x][\lforall[y][((x = y \land !A(x)) \lif !A(y))]]$
\item $\lexists[x][!A(x)] \land \lforall[y][\lforall[z][((!A(y) \land
    !A(z)) \lif y = z)]] \Sequent 
\lexists[x][(!A(x) \land \lforall[y][(!A(y) \lif y = x)])]$
\end{enumerate}
\end{prob}

\end{document}
